%
% Donné  en NFP 107 semestre 1 et en semestre 2 2023
%

\enteteExam{NFP 107 -- Systèmes de gestion de bases de données - Durée 2h30}%
{Examen FOAD session 2}%
{Documents autorisés\,: 10 pages recto-verso de notes manuscrites}%

Une application met en contact des personnes dont certaines (les professeurs)
proposent des leçons en ligne à d'autres (les étudiants). 
Dans ce qui suit on va supposer ques les leçons sont des cours de langue.


Au départ le concepteur a placé les données sur les leçons
dans un unique fichier CSV
dont voici un extrait du contenu sous forme de tableau.


\begin{center}
\begin{small}
\begin{tabular}{c|c|c|c|c|c}
\textbf{idProf} & \textbf{professeur} & \textbf{idEtudiant} &  \textbf{étudiant} & \textbf{langue} & \textbf{date} \\
\hline
\hline
1 & Luca & 2 & Philippe & italien & 14 mars\\
2 &  Philippe & 3 & Lydia & français & 15 mars\\
1 &  Luca & 3 & Lydia & italien & 20 mars\\
2 &  Philippe & 3 & Lydia & français & 2 avril \\
1 &  Luca & 3 & Lydia & italien & 6 avril\\
3 &  Lydia & 2 & Philippe & anglais & 12 avril \\
3 &  Lydia & 2 & Philippe & allemand & 13 avril \\
\hline
\end{tabular}
\end{small}
\end{center}

On va chercher à faire mieux avec une base relationnelle.

\subsection*{Conception du schéma (7 pts)}

\begin{enumerate}
   \item Pour chacune des DFs suivantes, indiquez si elle est respectée ou pas
   \emph{dans le tableau précédent}. Si elle ne l'est pas, expliquez pourquoi.
   
   \begin{itemize}
     \item $professeur \to langue$
     \item   $professeur, date \to etudiant$
     \item $etudiant \to professeur$
     \item $etudiant, langue \to professeur$
   \end{itemize}

   \item Finalement, après avoir étudié beaucoup plus d'exemples, on
      détermine les dépendances fonctionnelles suivantes: 
      $idProf \to professeur$, $idEtudiant \to etudiant$ et
      $(idProf, idEtudiant, date) \to langue$. 
      \begin{enumerate}
        \item Quelle est la clé de la relation\,?
        \item Quelle est la décomposition en 3FN\,?
      \end{enumerate}
      
   Pour la suite de l'examen nous travaillons avec le schéma suivant.
    
\begin{itemize}
\item Personne(\textbf{id}, nom, codeLangueNatale)
\item Langue (\textbf{code}, intitulé)
\item Cours(\textbf{id}, idProfesseur, idEtudiant, codeLangue, date)
\end{itemize}

  \item Enumérez les clés étrangères dans ce schéma
  \item Donnez le schéma entité-association correspondant
  \item Donnez les commandes \texttt{create table}
  \item Une des relations correspond-elle à une réification\,? Si oui
     indiquez ce que serait le schéma sans réification.
\end{enumerate}

   
\begin{solution}
\begin{enumerate}
   \item Réponses
   \begin{itemize}
     \item $professeur \to langue$\,: non, car on voit que Lydia enseigne l'anglais et l'allemand
     \item   $professeur, date \to etudiant$\,: oui, pas de contre-exemple dans le tableau (mais on en trouverait sans doute en continuant à accumuler les exemples)
     \item $etudiant \to professeur$, non, Philippe suit des cours avec Luca et Lydia
     \item $etudiant, langue \to professeur$\,: oui, pour une langue donnée, chaque étudiant n'a qu'un seul professeur
   \end{itemize}

   \item Réponses
      \begin{enumerate}
        \item Quelle est la clé de la relation\,? C'est le triplet  $(idProf, idEtudiant, date)$
        \item Quelle est la décomposition en 3FN\,? On applique la 
        normalisation et on trouve 3 relations\,: Professeur (id, nom),
     Etudiant (id, nom) et Cours(idProf, idEtudiant, date, langue). 
     Il est clair qu'il est préférable de fusionner les tables Professeur et Etudiant.
        \end{enumerate}

  \item Clés étrangères\,: codeLangueNatale, codeLangue,
      idProfesseur et idEtudiant dans Cours.
   \item Standard
    \item Standard
  \item Réification\,: la relation Cours peut être interprétée comme
      venant d'une relation ternaire (Prof, Etudiant, Date), ce qui
      donnerait une clé avec trois attributs sans réification (voir les DF ci-dessus).
      Après réification, on a simplifié la clé, mais il faut prendre 
      garde à ce qu'un professeur ne donne pas deux cours au même
      moment (avec une clause \texttt{unique}). 
\end{enumerate}

\end{solution}


\subsection*{SQL (7 pts)}

Sur le schéma donné précédemment, exprimez en SQL les requêtes suivantes\,:

\begin{enumerate}
  \item Liste des cours de l'étudiant Philippe (on suppose qu'il n'y en a qu'un), 
    en indiquant le nom
     du professeur et la date.
   \item Les personnes qui sont professeur \emph{et} étudiant
   \item Les professeurs qui donnent des cours dans une autre langue
       que leur langue natale
   \item Les professeurs qui ne donnent jamais cours dans une autre langue
       que leur langue natale
   \item Les personnes et les langues qu'elles n'apprennent pas
   \item Donner, pour chaque personne, le nombre de langues qu'elle apprend
   \item Donner les langues pour lesquelles on trouve au moins
       trois personnes dont c'est la langue natale. 
\end{enumerate}


\begin{solution}
\begin{enumerate}
  \item
\begin{verbatim}
select prof.nom as nomProfesseur, 
from Cours as c, Personne as p, Personne as prof
where c.idEtudiant = p.id
and   c.idProfesseur = prof.id
and    p.nom='Philippe'
\end{verbatim}

 \item Il faut trouver un cours où la personne est étudiante, un autre
   où la personne est professeure. Le distinct est utile.
\begin{verbatim}
select distinct p.nom 
from Cours as c1, Cours as c2, Personne as p
where c1.idEtudiant = p.id
and   c2.idProfesseur = p.id
\end{verbatim}


   \item il faut trouver (au moins) un cours d'une langue
       qui  n'est pas celle du professeur.
 \begin{verbatim}
select distinct p.nom 
from Cours as c, Personne as p
where c.codeLangue != p.codeLangueNatale
and   c.idProfesseur = p.id
\end{verbatim}

  \item Ici, il faut que l'ensemble des cours dans une autre langue que
    la langue natale du professeur soit vide.
    
\begin{verbatim}
select  p.nom
from Personne as p
where not exists (select *
    from Langue as l, Cours as c
    where   c.idProfesseur = p.id
    and c.codeLangue != p.codeLangueNatale)
\end{verbatim}

  \item On combine toutes les personnes et toutes les langues, 
  on ne garde que les paires pour lesquelles il n'existe pas de cours

\begin{verbatim}
select p.nom, l.intitulé
from Personne as p, Langue as l
where not exists (select*
    from Cours as c
    and   c.idEtudiant = p.id
    and c.codeLangue = l.code)
\end{verbatim}

    \item Petite astuce: il faut compter le nombre de langues distinctes
    (n'était pas pénalisé)
    
\begin{verbatim}
select p.nom, count(distinct langue.intitulé) as langue
from Personne as p, Cours as c, Langue as l
where c.idEtudiant = p.id
and c.codeLangue = l.code
group by p.id, p.nom 
\end{verbatim}
    
   \item 
   
\begin{verbatim}
select l.intitulé, count(*) as nbPersonnes
from Personne as p, Langue as l
where l.code = p.codeLangueNatale
group by l.code
having count(*)  >= 3
\end{verbatim}
\end{enumerate}

\end{solution}


\subsection*{Algèbre (2 pts)}


Donnez les expressions algébriques pour les requêtes suivantes\,:

\begin{enumerate}
  \item Nom des personnes qui se donnent un cours à elles-mêmes..
  \item Les personnes (on se contentera de l'identifiant)
    qui n'étudient pas l'italien (langue dont le code est \texttt{it})
\end{enumerate}

\begin{solution}
\begin{enumerate}
  \item $\pi_{nom} (\sigma_{idProfesseur=idEtudiant}(Cours) \Join_{idProfesseur=id} Personne)$
  \item $\pi_{id} (Personne) -  \pi_{idEtudiant} (\sigma_{langie='IT'} (Cours))$
 \end{enumerate}

\end{solution}


\subsection*{Valeurs nulles, vues (2 pts)}

\begin{enumerate}
   \item Créez une vue \texttt{Inscriptions} donnant, pour chaque
         personne, le code des langues apprises, ou \texttt{NULL}
         si aucune langue n'est apprise.
      
      \item En utilisant cette vue, quelle requête donne les personnes
        qui n'apprennent aucune langue\,?
\end{enumerate}

\begin{solution}
\begin{enumerate}
\item Il faut une jointure externe
\begin{verbatim}
create view Inscription as 
select p.nom, c.codeLangue
from Personne as p outer join Cours as c 
where c.idEtudiant = p.id
\end{verbatim}

\item 
\begin{verbatim}
select *
from Inscription as i
where i.codeLangue is null
\end{verbatim}
 \end{enumerate}

\end{solution}

\subsection*{Transactions (2 pts)}

Expliquez ce que signifie l'acronyme ACID.

\end{document}

